% Einheit 4 von 5 – Dezimalzahlen multiplizieren und dividieren (bank/bruchrechnung/e4.jsonl, zusammenbau v0.1)
%% TODO Zweigzeile Teil 1: was hier gelernt wird (Satz in Schülersprache aus dem Einheitstitel) – Einheit 4
\einheitenkopf[e4]{Einheit 4 von 5 · Dezimalzahlen multiplizieren und dividieren}
\zweigzeile{ab Kl. 5, je nach Buch bis Kl. 6 (am Gymnasium ab Kl. 6) · keine P10-Aufgabe}
\setcounter{aufgabe}{22}

%% TODO Titel: Ich-kann-Satz fehlt in der Bank (Platzhalter: Kettenname) – Nr. 23
%% TODO Anweisung: ein Satz über der ersten Teilaufgabe fehlt in der Bank – Nr. 23
\begin{aufgabe}{Wie viele Kommastellen hat das Ergebnis?}
\begin{teile}
\teil $3{,}2 \cdot 0{,}4$ – wie viele Kommastellen hat das Ergebnis? \leerfeld
\end{teile}
\end{aufgabe}

%% TODO Titel: Ich-kann-Satz fehlt in der Bank (Platzhalter: Kettenname) – Nr. 24
%% TODO Anweisung: ein Satz über der ersten Teilaufgabe fehlt in der Bank – Nr. 24
\begin{aufgabe}{Dezimal multiplizieren/dividieren}
%% TODO Darstellung für schwach fehlt in der Bank (bruchrechnung-e4-k2-s1-v1; Abschnitt „Für schwache Schüler“)
\swz{$3{,}72 \cdot 10$}{}
%% TODO Darstellung für schwach fehlt in der Bank (bruchrechnung-e4-k2-s1-v2; Abschnitt „Für schwache Schüler“)
\swz{$0{,}58 \cdot 100$}{}
%% TODO Darstellung für schwach fehlt in der Bank (bruchrechnung-e4-k2-s1-v3; Abschnitt „Für schwache Schüler“)
\swz{$4{,}5 \cdot 100$}{}
%% TODO Darstellung für schwach fehlt in der Bank (bruchrechnung-e4-k2-s1-v4; Abschnitt „Für schwache Schüler“)
\swz{$0{,}093 \cdot 10$}{}
%% TODO Darstellung für schwach fehlt in der Bank (bruchrechnung-e4-k2-s1-v5; Abschnitt „Für schwache Schüler“)
\swz{$12{,}6 \cdot 10$}{}
\swfrage{Was bleibt gleich, was ändert sich?}
%% TODO Darstellung für schwach fehlt in der Bank (bruchrechnung-e4-k2-s2-v1; Abschnitt „Für schwache Schüler“)
\swz{$45{,}7 : 10$}{}
%% TODO Darstellung für schwach fehlt in der Bank (bruchrechnung-e4-k2-s3-v1; Abschnitt „Für schwache Schüler“)
\swz{$2{,}3 \cdot 3$}{}
%% TODO Darstellung für schwach fehlt in der Bank (bruchrechnung-e4-k2-s4-v1; Abschnitt „Für schwache Schüler“)
\swz{$1{,}2 \cdot 0{,}7$}{}
%% TODO Darstellung für schwach fehlt in der Bank (bruchrechnung-e4-k2-s5-v1; Abschnitt „Für schwache Schüler“)
\swz{$0{,}2 \cdot 0{,}4$}{}
%% TODO Darstellung für schwach fehlt in der Bank (bruchrechnung-e4-k2-s6-v1; Abschnitt „Für schwache Schüler“)
\swz{$8{,}4 : 4$}{}
%% TODO Darstellung für schwach fehlt in der Bank (bruchrechnung-e4-k2-s7-v1; Abschnitt „Für schwache Schüler“)
\swz{$5{,}6 : 0{,}7$}{}
\swz[3]{Ein Lieferwagen braucht 8 l Diesel auf 100 km und fährt im Jahr $25\,000$ km. Ein Liter Diesel kostet $162{,}4$ ct. Berechne die Kraftstoffkosten für ein Jahr in Euro. \hfill (P10 2014 OS)}{}
\swz[3]{Welcher Wert ist am kleinsten: $2{,}2$; $0{,}22$; $0{,}2^2$; $22\,\%$? \hfill (P10 2023 OS)}{}
\end{aufgabe}

%% TODO Titel: Ich-kann-Satz fehlt in der Bank (Platzhalter: Kettenname) – Nr. 25
%% TODO Anweisung: ein Satz über der ersten Teilaufgabe fehlt in der Bank – Nr. 25
\begin{aufgabe}{Dezimal multiplizieren/dividieren – Fehler finden, Begründen, Anwendung}
\swz[3]{Kim rechnet: $0{,}7 \cdot 0{,}4 = 2{,}8$. Finde den Fehler.}{}
\swfrage{Warum rückt das Komma bei $\cdot 10$ eine Stelle nach rechts?}
\swz[3]{Äpfel kosten lose 2,40 € je kg. Ein Beutel mit 1,5 kg kostet 3,29 €. Was ist günstiger?}{}
\end{aufgabe}

\uebersichtskasten{Mal 10, 100, 1000: Komma nach rechts; geteilt durch 10, 100: Komma nach links. \quad 0,46 · 100 = 46 \quad 3,8 : 10 = 0,38 \\ Mal: ohne Komma rechnen, dann Kommastellen beider Zahlen zählen. \quad 1,5 · 0,04 = 0,06 (15 · 4 = 60, drei Kommastellen) \\ Geteilt: beim Teiler das Komma wegschieben, beim anderen genauso weit. \quad 7,2 : 0,9 = 72 : 9 = 8}
